% ============================================================
% glossaire_termes.tex
% Définitions des termes et expressions utilisés dans le rapport
% (distinct de glossaire.tex, qui ne porte que les acronymes)
% NOTE : une entrée n'apparaît dans le glossaire que si elle est
% appelée au moins une fois via \gls{clé} dans un chapitre, sauf
% usage de \glsaddall (voir front/glossaire.tex).
% ============================================================

\newglossaryentry{serveurphysique}{
  name={Serveur},
  description={Machine hôte, branchée à l'outil, hébergeant
    des conteneurs et exposant ses ressources à la supervision}
}

\newglossaryentry{conteneur}{
  name={Conteneur},
  description={Unité d'exécution isolée hébergeant une application,
    rattachée à un serveur hôte et à un environnement}
}

\newglossaryentry{application}{
  name={Application},
  description={Logiciel métier exécuté dans un ou plusieurs conteneurs}
}

\newglossaryentry{depot}{
  name={Dépôt},
  description={Source de code versionnée, indépendante de la ressource
    physique sur laquelle elle est déployée}
}

\newglossaryentry{deploiement}{
  name={Déploiement},
  description={Mise à disposition d'une application à partir d'un
    dépôt, activable ou désactivable par serveur et par dépôt}
}

\newglossaryentry{sauvegarde}{
  name={Sauvegarde},
  description={Copie horodatée, restaurable et vérifiée d'une base de
    données, activable ou désactivable par serveur}
}

\newglossaryentry{metrique}{
  name={Métrique},
  description={Mesure de consommation d'une ressource (processeur,
    mémoire, disque) relevée à la demande ou en continu}
}

\newglossaryentry{seuil}{
  name={Seuil},
  description={Valeur limite définie pour une ressource, servant de
    référence au déclenchement d'une alerte}
}

\newglossaryentry{alerte}{
  name={Alerte},
  description={Signalement déclenché au franchissement d'un seuil}
}

\newglossaryentry{supervision}{
  name={Supervision},
  description={Surveillance centralisée de l'état et des métriques de
    l'ensemble du parc, à la demande ou en continu}
}

\newglossaryentry{consoleweb}{
  name={Console web},
  description={Moyen d'exécuter des commandes sur un serveur ou un
    conteneur directement depuis l'interface, dont chaque session est
    intégralement journalisée}
}

\newglossaryentry{perimetre}{
  name={Périmètre (scope)},
  description={Étendue sur laquelle porte une permission : globale, ou
    restreinte à un serveur, un conteneur ou un dépôt précis}
}

\newglossaryentry{permission}{
  name={Permission},
  description={Droit d'effectuer une action donnée, dans un périmètre
    défini}
}

\newglossaryentry{secret}{
  name={Secret},
  description={Donnée sensible propre à une ressource (clé SSH, jeton),
    chiffrée et non restituée en clair par l'API}
}
